% 【20260901 重写】原模板逐字教了评委连续点名的两个毛病，改的人照模板写，所以四轮没改掉：
%   ① `\textbf{对于问题一,}` —— 审稿员 2A/2B/3A 连续三轮抱怨「摘要连续使用『对于问题一/二/三』」
%   ② `首先…接着…最后…` —— 运行时/审计.py 的「首先最后套件」判据正在罚这个句式
%   ③ 「各模型性能优越、结论可靠」「提供理论依据和决策参考」—— 典型 AI 论文套话
% 下面按 运行时/优秀论文标准.md 第 3、5 条重写。**方括号内是写法要求，不是要填的占位符。**

\begin{abstract}

% ==== 开头（1-2 句）：交代这道题的物理/业务实质，不要「…是…领域的重要研究课题」 ====
% [写清研究对象的关键特征与本文据以立论的数据，让评委立刻知道这是哪一类问题]

% ==== 逐问（每问一段）====
% 段首用**问题特征**开头，不用「对于问题N」。审稿员给的范例：
%   「在一次往返边界下…」「对碳化硅双角谱…」「面对硅片的高阶往返…」
% 每段的顺序固定为：先朴素回答题目问的那个量 → 再说凭什么信 → 最后才是方法名。
% 反例（评委原话点名）：「问题三 Airy 形态正检出但机制未独立确认，厚度不可交付。」
% 正例：「硅片光谱确实出现了多光束干涉的形态特征，但支持这一机制的独立证据不足，
%        因此本文不据此给出硅片厚度。」
% 方法名每问只留 1 个必要的，其余缩写移到正文；术语首次出现补一句白话定义。

% ==== 结尾（1-2 句）====
% 写**本文结论的适用边界与它依赖什么**，不要「性能优越、结论可靠」「提供理论依据和决策参考」。
% 例：「上述结论建立在 X 的前提上；若 Y 可实测，Z 的不确定度可由 A 收窄到 B。」

\keywords{关键词1；关键词2；关键词3；关键词4；关键词5}

\end{abstract}
